\chapter*{Glossaire}
\addcontentsline{toc}{chapter}{Glossaire}

\vspace{0.5cm}

Ce glossaire regroupe les principaux termes techniques et acronymes 
utilisés dans ce rapport de stage.

\vspace{0.5cm}

\begin{description}[leftmargin=4cm, labelwidth=3.5cm]

    \item[\textbf{API}] \textit{Application Programming Interface} — 
    Interface de programmation permettant à deux applications de 
    communiquer entre elles. Dans ce projet, les APIs REST permettent 
    la communication entre la WebView, le dashboard monitor et le 
    back-end Spring Boot.

    \vspace{0.3cm}

    \item[\textbf{Agile}] Ensemble de méthodologies de gestion de 
    projet basées sur le développement itératif et incrémental, 
    la collaboration et l'adaptation au changement. Le Manifeste 
    Agile, publié en 2001, définit les valeurs et principes fondateurs 
    de ces méthodologies.

    \vspace{0.3cm}

    \item[\textbf{Back-end}] Partie serveur d'une application web 
    ou mobile, invisible pour l'utilisateur final. Le back-end 
    gère la logique métier, le traitement des données et la 
    communication avec la base de données.

    \vspace{0.3cm}

    \item[\textbf{Dashboard}] Tableau de bord interactif permettant 
    de visualiser, analyser et gérer des données de manière centralisée. 
    Dans ce projet, le dashboard monitor permet à l'équipe Bewize 
    de consulter et traiter les feedbacks et réclamations reçus.

    \vspace{0.3cm}

    \item[\textbf{Docker}] Plateforme de conteneurisation permettant 
    d'empaqueter une application et ses dépendances dans des 
    conteneurs isolés et portables, garantissant un comportement 
    identique dans tous les environnements.

    \vspace{0.3cm}

    \item[\textbf{docker-compose}] Outil Docker permettant de définir 
    et d'orchestrer plusieurs conteneurs Docker via un fichier de 
    configuration YAML. Utilisé dans ce projet pour orchestrer 
    les conteneurs Spring Boot, PostgreSQL et React JS.

    \vspace{0.3cm}

    \item[\textbf{Feedback}] Retour d'information d'un utilisateur 
    sur un produit ou service. Dans ce projet, le feedback désigne 
    un retour positif ou constructif soumis par un utilisateur 
    de l'application mobile Bewize.

    \vspace{0.3cm}

    \item[\textbf{Front-end}] Partie cliente d'une application web 
    ou mobile, visible et accessible par l'utilisateur final. 
    Le front-end gère l'interface utilisateur et l'expérience 
    visuelle de l'application.

    \vspace{0.3cm}

    \item[\textbf{Git}] Système de contrôle de version distribué 
    permettant de suivre les modifications du code source, de 
    collaborer en équipe et de gérer différentes versions d'un projet.

    \vspace{0.3cm}

    \item[\textbf{GitHub}] Plateforme web basée sur Git offrant 
    un hébergement de dépôts Git en ligne et des fonctionnalités 
    collaboratives avancées (pull requests, issues, CI/CD).

    \vspace{0.3cm}

    \item[\textbf{HTTP}] \textit{HyperText Transfer Protocol} — 
    Protocole de communication utilisé pour le transfert de données 
    sur le web. Les APIs REST de ce projet utilisent HTTP pour 
    la communication entre les composants.

    \vspace{0.3cm}

    \item[\textbf{JPA}] \textit{Java Persistence API} — Spécification 
    Java pour la gestion de la persistance des données et le mapping 
    objet-relationnel (ORM). Utilisée via Spring Data JPA pour 
    la communication avec PostgreSQL.

    \vspace{0.3cm}

    \item[\textbf{JSON}] \textit{JavaScript Object Notation} — 
    Format léger d'échange de données, facile à lire pour les 
    humains et à analyser pour les machines. Utilisé pour les 
    échanges de données entre la WebView et l'API Spring Boot.

    \vspace{0.3cm}

    \item[\textbf{Node.js}] Environnement d'exécution JavaScript 
    côté serveur, basé sur le moteur V8 de Chrome. Utilisé comme 
    environnement d'exécution pour React JS et pour la gestion 
    des dépendances via npm.

    \vspace{0.3cm}

    \item[\textbf{npm}] \textit{Node Package Manager} — Gestionnaire 
    de paquets officiel de Node.js, permettant d'installer et de 
    gérer les bibliothèques JavaScript nécessaires au projet.

    \vspace{0.3cm}

    \item[\textbf{PostgreSQL}] Système de gestion de base de données 
    relationnelle open-source, reconnu pour sa robustesse et sa 
    conformité aux standards SQL. Utilisé pour stocker les feedbacks 
    et réclamations dans ce projet.

    \vspace{0.3cm}

    \item[\textbf{Product Backlog}] Liste ordonnée de toutes les 
    fonctionnalités à développer pour le produit, gérée par le 
    Product Owner dans le cadre de la méthodologie Scrum.

    \vspace{0.3cm}

    \item[\textbf{React JS}] Bibliothèque JavaScript open-source 
    développée par Meta pour construire des interfaces utilisateur 
    web interactives et réactives, basée sur le concept de composants 
    réutilisables.

    \vspace{0.3cm}

    \item[\textbf{React Native}] Framework open-source développé 
    par Meta permettant de créer des applications mobiles natives 
    pour iOS et Android à partir d'un seul code source JavaScript.

    \vspace{0.3cm}

    \item[\textbf{Réclamation}] Retour négatif ou signalement d'un 
    problème soumis par un utilisateur de l'application mobile 
    Bewize, nécessitant une action corrective de la part de l'équipe.

    \vspace{0.3cm}

    \item[\textbf{REST}] \textit{Representational State Transfer} — 
    Style d'architecture pour les APIs web, basé sur les protocoles 
    HTTP standard. Les APIs développées dans ce projet suivent 
    les principes REST.

    \vspace{0.3cm}

    \item[\textbf{Scrum}] Framework agile de gestion de projet 
    organisé en Sprints itératifs, avec trois rôles (Product Owner, 
    Scrum Master, équipe de développement) et trois artefacts 
    (Product Backlog, Sprint Backlog, incrément).

    \vspace{0.3cm}

    \item[\textbf{Sprint}] Itération de durée fixe (généralement 
    une à quatre semaines) dans le framework Scrum, à l'issue 
    de laquelle l'équipe livre un incrément fonctionnel du produit.

    \vspace{0.3cm}

    \item[\textbf{Sprint Backlog}] Liste des tâches sélectionnées 
    pour un Sprint donné, créée lors du Sprint Planning et gérée 
    par l'équipe de développement.

    \vspace{0.3cm}

    \item[\textbf{Spring Boot}] Framework Java open-source simplifiant 
    le développement d'applications back-end et de microservices, 
    basé sur le principe de convention over configuration.

    \vspace{0.3cm}

    \item[\textbf{UML}] \textit{Unified Modeling Language} — Langage 
    de modélisation standardisé permettant de représenter graphiquement 
    la structure et le comportement d'un système logiciel.

    \vspace{0.3cm}

    \item[\textbf{User Story}] Description courte d'une fonctionnalité 
    exprimée du point de vue de l'utilisateur final, suivant le 
    format : "En tant que [rôle], je veux [fonctionnalité] afin 
    de [bénéfice]".

    \vspace{0.3cm}

    \item[\textbf{WebView}] Composant permettant d'afficher une 
    page web directement dans une application mobile native, sans 
    rediriger l'utilisateur vers un navigateur externe.

\end{description}